\documentclass{report}
\usepackage{amsmath,amssymb}
\setlength{\parindent}{0mm}
\setlength{\parskip}{1em}
\begin{document}
\begin{center}
\rule{6in}{1pt} \
{\large Yongxiang Liu \\
{\bf A Dirichlet-Neumann Method for Plate Problem}}

Section of Mathematics \\ University of Geneva \\ Rue du Lievre 2-4 \\ CP 64 \\ 1211 Geneva 4 \\ Switzerland
\\
{\tt Yongxiang.Liu@unige.ch}\\
Martin Gander\end{center}

We study a Dirichlet-Neumann method for the plate bending problem. This
is a forth order problem and thus requires four interface conditions to
hold between subdomains, which is quite different from the classical
Laplace problem which requires only the continuity of the Dirichlet and
Neumann traces at interfaces. We present a particular Dirichlet-Neumann
iterative method for a specific choice
of interface conditions and a suitably chosen relaxation matrix. We
determine the optimal choice of the relaxation matrix, and prove that the
associated convergence rate is then independent of the mesh size. We
illustrate our analysis with numerical experiments.


\end{document}
